%%% FIELD prop-statement
Uniformly over \(2\le d\le n\) and \(0\le\sigma\le2\) satisfying
\[
 \sigma^2d^2\le n
 \quad\text{or}\quad
 u_{n,d}\le \frac1n,
\]
equivalently satisfying the first inequality or \(d^2\le n\log^2(en)\),
\[
 \mathsf R_{n,d,\epsilon,M,\sigma}\asymp_{\epsilon}\frac{M^2}{n}.
\]
In particular, for \(d=\lfloor n^{3/4}\rfloor\) and \(\sigma^2=u_{n,d}\), the minimax mean squared error is of order \(M^2/n\).
%%% FIELD prop-proof
Fix admissible \(n,d,\sigma\) with \(2\le d\le n\).  The lower inequality of \(\mathrm{thm:refined\text{-}bracket}\), followed by the definition \(b_{n,d}=n^{-1}+dn^{-2}\), gives
\[
 \mathsf R_{n,d,\epsilon,M,\sigma}
 \ge c_{\epsilon}M^2\{b_{n,d}+\sigma^2\min(1,u_{n,d})\}
 \ge c_{\epsilon}M^2b_{n,d}
 \ge c_{\epsilon}\frac{M^2}{n}.                                      \tag{1}
\]
This lower bound applies in both branches of the asserted union.

Suppose first that \(\sigma^2d^2\le n\).  Because \(d\le n\), \(\mathrm{lem:empirical\text{-}mass\text{-}upper}\) applies to the total estimator \(\widehat\tau^{\mathrm{em}}\).  The definition of minimax risk and that lemma imply
\[
 \begin{aligned}
 \mathsf R_{n,d,\epsilon,M,\sigma}
 &\le \sup_{P\in\mathcal P_{n,d,\epsilon,M,\sigma}}
 E_{Q_P^{(n)}}[(\widehat\tau^{\mathrm{em}}-\tau(P))^2]\\
 &\le C_{\epsilon}M^2\left(\frac1n+\sigma^2\frac{d^2}{n^2}\right)
 \le 2C_{\epsilon}\frac{M^2}{n}.                                    \tag{2}
 \end{aligned}
\]
The last inequality uses exactly \(\sigma^2d^2\le n\).

Suppose instead that \(u_{n,d}\le n^{-1}\).  By the upper inequality of \(\mathrm{thm:refined\text{-}bracket}\) and the definition of \(q_{n,d,\sigma}^{\mathrm{em}}\),
\[
 \begin{aligned}
 \mathsf R_{n,d,\epsilon,M,\sigma}
 &\le C_{\epsilon}M^2q_{n,d,\sigma}^{\mathrm{em}}\\
 &=C_{\epsilon}M^2\left[\frac1n+
 \min\left\{1,u_{n,d},\frac d{n^2}+\sigma^2v_{n,d}\right\}\right]\\
 &\le C_{\epsilon}M^2\left(\frac1n+u_{n,d}\right)
 \le 2C_{\epsilon}\frac{M^2}{n}.                                    \tag{3}
 \end{aligned}
\]
Combining (1) with (2) or (3) proves the uniform comparison.  Since \(n,d\), and \(\log(en)\) are positive,
\[
 u_{n,d}\le\frac1n
 \quad\Longleftrightarrow\quad
 \frac{d^2}{n^2\log^2(en)}\le\frac1n
 \quad\Longleftrightarrow\quad
 d^2\le n\log^2(en),                                                   \tag{4}
\]
which proves the stated equivalent form.

Finally take \(d=\lfloor n^{3/4}\rfloor\) and \(\sigma^2=u_{n,d}\).  For every \(n\ge3\), \(2\le d\le n\), and
\[
 \sigma^2d^2
 =\frac{d^4}{n^2\log^2(en)}
 \le\frac{n}{\log^2(en)}
 \le n.                                                               \tag{5}
\]
Thus this sequence is in the first branch, and (1)--(2) give the asserted order \(M^2/n\).
%%% FIELD residual-construction
\[
 \mathcal W:=\left\{(n,d,\sigma):
 n\ge3,\quad
 2\le d\le\left\lfloor\frac{n^2}{\log(en)}\right\rfloor,\quad
 0<\sigma<2,\quad
 \frac1n<u_{n,d}<1,\quad
 \bigl[d>n\ \text{or}\ \sigma^2d^2>n\bigr]
 \right\}.
\]
%%% FIELD oeq-statement
At \(d=n\) and \(\sigma=\sigma_n^\star\), carry out \(\mathrm{def:benchmark\text{-}repair\text{-}handle}\): determine the maximal \(\rho_n\) for supportwise feasible, exactly normalized priors with identical design marginals, ATEs concentrated in disjoint intervals separated by \(M\rho_n\), and full-sample total variation at most \(1/3\). Prove either \(\rho_n\sqrt n\to\infty\), which supplies a substantive same-class lower bound inside the accepted parent's residual wedge, or a quantitative obstruction showing that this conditionally centered shared-design program cannot exceed the parametric separation order. Only a benchmark construction that passes exact normalization, realization-wise radius, target concentration, and full-mixture distance may be extended across \(\mathcal W\); no global rate formula is committed before that test.
%%% FIELD oeq-partial
Define \(\rho_n^{\mathrm{outer}}\) as the supremum of the normalized distances between the two target-support intervals over all pairs of priors satisfying only exact normalization, supportwise membership in \(\mathcal P_{n,n,\epsilon,M,\sigma_n^\star}\), identical design marginals, and full-sample total variation at most \(1/3\), without imposing a particular centered-score representation.  The frozen results imply the rigorous outer bracket
\[
 c n^{-1/2}\le \rho_n^{\mathrm{outer}}
 \le C_{\epsilon}\log(en)n^{-1/2}.                                   \tag{6}
\]
Here is the derivation.  At \(d=n\) one has \(v_{n,n}=1\), \(u_{n,n}=\log^{-2}(en)\), and \((\sigma_n^\star)^2=\log^2(en)/n\).  The third branch of \(q_{n,n,\sigma_n^\star}^{\mathrm{em}}\) and \(\mathrm{thm:refined\text{-}bracket}\) therefore give
\[
 \begin{aligned}
 \mathsf R_{n,n,\epsilon,M,\sigma_n^\star}
 &\le C_{\epsilon}M^2\left\{\frac1n+\frac1n+
 (\sigma_n^\star)^2\right\}\\
 &\le 3C_{\epsilon}M^2(\sigma_n^\star)^2.                            \tag{7}
 \end{aligned}
\]
The last step uses \(\log(en)>1\).

Consider any feasible pair of priors whose target-support intervals \(I_{-}\) and \(I_{+}\) have distance at least \(M\rho\), and let their sample-mixture laws be \(\mathbb Q_{-}\) and \(\mathbb Q_{+}\).  Given any estimator \(T\), classify its value according to the nearer of \(I_{-}\) and \(I_{+}\), breaking ties deterministically.  If this classifier errs under hypothesis \(h\), then
\[
 |T-\tau(P)|\ge \operatorname{dist}(T,I_h)\ge \frac{M\rho}{2}.         \tag{8}
\]
Writing \(e_h\) for its error probability under \(\mathbb Q_h\), averaging (8) over the two priors yields
\[
 \frac12\sum_{h\in\{-1,1\}}\int
 E_{Q_P^{(n)}}[(T-\tau(P))^2],d\Pi_h(P)
 \ge \frac{M^2\rho^2}{8}(e_{-}+e_{+}).                               \tag{9}
\]
The elementary testing identity \(e_{-}+e_{+}\ge1-\operatorname{TV}(\mathbb Q_{-},\mathbb Q_{+})\), obtained by writing the acceptance region of the classifier in the variational definition of total variation, and feasibility give
\[
 \sup_{P\in\mathcal P_{n,n,\epsilon,M,\sigma_n^\star}}
 E_{Q_P^{(n)}}[(T-\tau(P))^2]
 \ge \frac{M^2\rho^2}{12}.                                           \tag{10}
\]
Taking the infimum over \(T\), comparing (10) with (7), and taking square roots proves
\[
 \rho\le 6\sqrt{C_{\epsilon}}\,\sigma_n^\star
 =6\sqrt{C_{\epsilon}}\,\frac{\log(en)}{\sqrt n}.                  \tag{11}
\]

For the lower side of (6), fix \(0<c\le1/3\).  Under hypothesis \(h\in\{-1,1\}\), use the singleton prior on the law with \(p_k=1/n\), \(\pi_k=1/2\), and \(\tau_k=hcM/\sqrt n\) for every \(k\).  Set \(\mu_{ak}=(2a-1)\tau_k/2\).  Conditional on \(X=k\), draw \(A\sim\operatorname{Bernoulli}(1/2)\), draw potential outcomes independently of \(A\) with
\[
 P\{Y(a)=M/2\mid X=k\}
 =\frac12+\frac{(2a-1)\tau_k}{2M},                                   \tag{12}
\]
and set \(Y=Y(A)\).  Since \(|\tau_k|=cM/\sqrt n\le M\), (12) is a probability.  It gives
\[
 E[Y(a)\mid X=k]=\frac{(2a-1)\tau_k}{2}=\mu_{ak},
 \qquad
 \operatorname{Var}(Y(a)\mid X=k)
 =\frac{M^2}{4}-\mu_{ak}^2\le M^2.                                  \tag{13}
\]
Thus consistency follows from \(Y=Y(A)\), exchangeability follows because \(A\) is drawn independently of the potential outcomes conditionally on \(X\), overlap holds because \(\pi_k=1/2\in[\epsilon,1-\epsilon]\), the mean and variance envelopes follow from (13), and the radius is zero because every \(\tau_k\) equals \(\tau(P)=hcM/\sqrt n\).  Hence both laws belong to the frozen class.  Their ATE singleton supports are separated by \(2cM/\sqrt n\).

Put \(t=c/(2\sqrt n)\).  Conditional on either value of \((X,A)\), the two hypotheses induce Bernoulli probabilities \(1/2+t\) and \(1/2-t\) on the sign of \(Y\), in one order or the other.  Their Hellinger affinity is
\[
 2\sqrt{(1/2+t)(1/2-t)}=\sqrt{1-4t^2}.                                \tag{14}
\]
The design law is identical, so (14) is also the one-observation affinity.  Independence makes the \(n\)-observation affinity
\[
 \mathcal A_n=(1-4t^2)^{n/2}.                                         \tag{15}
\]
For dominated laws with densities \(p,q\), Cauchy--Schwarz applied to \(|p-q|=|\sqrt p-\sqrt q|(\sqrt p+\sqrt q)\) gives \(\operatorname{TV}(p,q)\le\sqrt{1-(\int\sqrt{pq})^2}\).  Using (15) and Bernoulli's inequality, which applies because \(0\le4t^2\le1\), yields
\[
 \operatorname{TV}(\mathbb Q_{-},\mathbb Q_{+})
 \le\sqrt{1-(1-4t^2)^n}
 \le2\sqrt n\,t=c\le\frac13.                                      \tag{16}
\]
This feasible pair has \(\rho=2c/sqrt n\), proving the lower side of (6).  The pair is a small global shift and does not instantiate the undefined conditionally centered marked-score mechanism.  Thus (6) neither produces a super-parametric feasible centered coupling nor proves a parametric obstruction for every such coupling.
